\section*{Closing Group Reflection}
\addcontentsline{toc}{section}{Closing Group Reflection}
\label{sec:reflection}

AI was most useful at the start of each task. It gave a fast overview of the actors, the sides of the platform and which partner controls which part of the infrastructure. Most basic facts were correct: Thales controls S3NS, Google supplies the technology, and ANSSI qualifies the service. In Tasks 3 and 4 it also invented no prices, probabilities or costs.

The problems were in the economic reasoning. In Task 1 the AI described the dependency between S3NS and Google as balanced, but the case sources show that S3NS needs Google far more than the reverse. In Task 2 it jumped from high switching costs to a winner-takes-all market, although keeping clients is not the same as winning the whole market. In Task 3 it listed Google and Thales as sides of the platform, but sides are user groups that interact through the platform, so we kept only clients and complementors. In Task 4 it said the cheapest partner should invest, which only works for total effort, while most of PREMI3NS is weakest link.

These errors share a pattern. The text sounded confident and used terms such as ``network effects'' and ``subsidy'' loosely. It mixed up economies of scale with network effects, and suppliers with sides. We only saw this because we had the course definitions and the case sources next to us.

This shows the limits of AI. It can collect and structure facts, but it does not decide which concept applies, or warn us when its answer is too strong. The main tension in this case, that S3NS sells sovereignty but depends on Google, came from our own analysis.
